\section{Conclusions}
\label{sec:conclusions}

\begin{enumerate}[leftmargin=1.6em,itemsep=4pt,topsep=4pt,parsep=0pt]

\item The frequency range is what makes this experiment new: no published
radio technosignature search above about 30\,GHz has reached a star within
40\,pc (\S\ref{sec:background}), and an archive assembled for other science
is at present the only way to reach one.

\item The primary experiment is \NWinA{} fine-channel windows toward
\NSysClassA{} stellar systems, covering \CvUnionA\,GHz of union bandwidth in
disconnected intervals, searched for an unresolved carrier whose frequency
drifts. A supplementary coarse-channel search for an unresolved spectral
excess adds \NWinB{} windows toward \PrSysClassB{} systems and
\CvUnionB\,GHz of union bandwidth, with no drift discrimination. Both come
from one uniform reduction of \NEB{} public execution blocks,
\TotalDataTB\,TB of raw data, and between them they span
\SurvFreqLo--\SurvFreqHi\,GHz.

\item Sensitivity is a measured completeness, not a threshold. Carriers
injected into the real calibrated visibilities are recovered through the
unchanged selection --- the trigger, and localisation at the star --- and
$\mathrm{EIRP}_{90}$ is the power at which nine in ten come back. The median
system's best window reaches $\EirpNinetySysMedA$\,W, with a spread across
systems of $\EirpNinetySysLoA$--$\EirpNinetySysHiA$\,W. The rank and the
visibility fit decide nothing, so neither is charged to that limit. One bias
runs one way only and is not inside the quoted calibration: an injected
carrier suffers no atmospheric decorrelation where a real one would, so
every limit here is optimistic by up to \SensDecorPct{} per cent.

\item All \EvNCrossRaw{} threshold crossings are reported with their
dispositions, \PrFlagXrossWord{} of them in the one execution block
that fails the data-quality criterion of
Appendix~\ref{sec:blockquality} and reported apart from the
\EvNCrossQp{} in quality-passing data; of those, \EvNAttrQp{} coincide with
a catalogued molecular transition, almost all of them carbon monoxide ---
including the $\beta$~Pictoris disc, which the search recovered without
being told where to look --- and \EvNUnattrQp{} do not.

\item Every unattributed crossing was tested for recurrence in the star's
own frame, in each further execution block covering the event frequency and
over the whole of each such window at every drift searched, because a
carrier on an accelerating platform need not return to one channel. None is
present again: the largest statistic at any predicted cell is \RcTMax{}
against $T_{\rm trig}=\StageNullTrig$, and the exceedances found anywhere in
those windows are no more than chance predicts for trials of that size.
Where a discovery amplitude survives, the repeat coverage rules out a
carrier of that amplitude in \RcNExcl{} cases, and one of half that
amplitude or less in \RcNFracHalf{} of them; taken instead as a
disagreement between two measurements of the same amplitude, discovery and
repeat differ by \RcExclMin{} to \RcExclMax$\sigma$, median
\RcExclMed$\sigma$, which is as much as the pair can say when the discovery
amplitude is itself a five-sigma measurement. Nor did \PrRankClause{}
return. What these tests exclude is a carrier that holds its frequency
inside the window searched and radiates through both observations, and not a
persistent transmitter of any kind.

\item What this survey could search and what it could confirm are different
extents. The carrier search covered
\RcSearchSys{} systems, of which only \RcConfSysDay{} --- carrying
\CvConfPctDay{} per cent of the searched union bandwidth --- were observed
again at the same frequency more than a day apart, and \RcConfSysYr{} beyond
a year. Across both classes, \RcNoRepSysAll{} of the \RcSearchSysAll{}
systems have no second block covering a frequency the first one covered, and
toward those the null reported here says nothing at any strength.

\item Combining each surveyed star's epochs in its own rest frame deepens
the flux limit by a median factor \SyStkGainMed. For Proxima~Centauri,
where $\sqrt{N_{\rm eff}}=\SxProxSqrtNeff$ is what combining the epochs
buys, the stacked limit measured by injection into the stack reaches
$\SzProxEirpCensus$\,W over the primary census alone and
$\StkProxEirpCal$\,W once the rest of the archive is added, in
\ChanBModeMHz\,MHz channels --- a Class~B limit on an unresolved excess,
carrying no drift discrimination --- for an emitter steady in the stellar
frame to within one channel, \StkProxTolKms\,km\,s$^{-1}$ peak to peak,
across \StkProxSpanYr\,yr.

\item These results constrain unresolved carriers only within the domain
actually searched and only toward this archive-selected sample
(\S\ref{sec:excludes}). What transfers beyond it is the obstacle: a
channelisation chosen by other observers for other purposes, and not
collecting area, is what decides which kind of search is possible at all.

\end{enumerate}
